\begin{figure}[tbp]
\centering
\setlength{\fboxsep}{7pt}
\fbox{\begin{minipage}{0.90\linewidth}
\centering
\textbf{One integral coefficient polynomial}\\[3pt]
$\displaystyle G_D(X)=[s^qt^q]\per_k(X+sI+tD)
\in\mathbb Z[\theta_1,\ldots,\theta_k][X_{ij}]$\\[3pt]
$\displaystyle D=\diag(\theta_1,\ldots,\theta_k),\qquad
r=k-2q,\qquad T=\binom{q+2}{2}$
\end{minipage}}

\smallskip
\begin{tabular}{@{}c@{\hspace{0.025\linewidth}}c@{}}
\begin{minipage}{0.43\linewidth}\centering
$\swarrow$\quad over $\C$, every parameter choice
\end{minipage}
&
\begin{minipage}{0.43\linewidth}\centering
$\searrow$\quad over $K=\overline{\mathbb F}_2$, choose distinct parameters
\end{minipage}
\\[4pt]
\fbox{\begin{minipage}{0.405\linewidth}\centering
\textbf{Permanent projection}\\[3pt]
$\displaystyle \per_m(A_m(X))=a_0G_D(X)$\\[3pt]
$m\ge11k-2q,\qquad a_0\ne0$
\end{minipage}}
&
\fbox{\begin{minipage}{0.405\linewidth}\centering
\textbf{Plane-curve geometry}\\[3pt]
$\displaystyle \nabla G_{\bar D}(X)=0
\ \Longrightarrow\ \delta(C_X)\ge T$\\[3pt]
$\displaystyle \codim V(\nabla G_{\bar D})\ge T-k$
\end{minipage}}
\\[5pt]
\begin{minipage}{0.43\linewidth}\centering
$\searrow$\quad apply at the complex parameters supplied by geometry
\end{minipage}
&
\begin{minipage}{0.43\linewidth}\centering
$\swarrow$\quad proper specialization and a transverse linear section
\end{minipage}
\end{tabular}

\smallskip
\fbox{\begin{minipage}{0.90\linewidth}
\centering
\textbf{A smooth affine restriction over $\C$}\\[3pt]
A nonzero homogeneous form $G$ of degree $r$ in $d$ variables has a smooth
projective zero set and satisfies $G=\per_m\circ\Phi_m$ for a fixed affine map $\Phi_m$.\\[3pt]
$\displaystyle r-1\ge\frac{k}{4},\qquad
d-1\ge\frac{k^2}{200}$
\end{minipage}}

\smallskip
$\big\downarrow$\quad the polar count for smooth homogeneous forms
\smallskip

\fbox{\begin{minipage}{0.90\linewidth}
\centering
$\displaystyle
\bdc(\per_m)\ge\bdc(G)
\ge\frac{(r-1)(d-1)}{4e}
\ge\frac{k^3}{3200e}$
\end{minipage}}
\caption{The two uses of the integral coefficient formula meet over
$\C$. The graph projection applies to every complex parameter choice;
the characteristic-two branch supplies a suitable complex choice through
the proper parameter family. The displayed scale bounds use
$q=\lfloor(k-1)/3\rfloor$, $k\ge128$, and $m\ge11k-2q$.}
\label{fig:proof-flow}
\end{figure}
